% Estilo común de los manuales del gemelo digital (se incluye con \input)
\usepackage{geometry}
\geometry{letterpaper, margin=2.2cm}
\usepackage[spanish, es-nodecimaldot, mexico]{babel}
\usepackage[utf8]{inputenc}
\usepackage[T1]{fontenc}
\usepackage{lmodern}
\usepackage{amsmath,amssymb}
\usepackage[table]{xcolor}
\definecolor{naranja}{RGB}{194,98,12}
\definecolor{rojo}{RGB}{192,57,43}
\definecolor{verde}{RGB}{30,120,70}
\definecolor{azulw}{RGB}{0,103,184}
\definecolor{azulg}{RGB}{70,90,120}
\definecolor{gris}{RGB}{242,242,242}
\definecolor{grisos}{RGB}{90,90,90}
\usepackage{booktabs,longtable,array,tabularx}
\usepackage{graphicx}
\usepackage{float}
\usepackage{caption}
\captionsetup{font=small, labelfont=bf}
\usepackage{tikz}
\usetikzlibrary{positioning,arrows.meta,fit,calc,shapes.geometric,decorations.pathreplacing,backgrounds}
\usepackage{enumitem}
\setlist{itemsep=2pt}
\usepackage{listings}
\lstset{basicstyle=\ttfamily\small, frame=single, backgroundcolor=\color{gris}, breaklines=true,
  columns=fullflexible, keepspaces=true, extendedchars=true, upquote=true,
  literate={á}{{\'a}}1 {é}{{\'e}}1 {í}{{\'i}}1 {ó}{{\'o}}1 {ú}{{\'u}}1 {ñ}{{\~n}}1 {Á}{{\'A}}1 {É}{{\'E}}1
           {Í}{{\'I}}1 {Ó}{{\'O}}1 {Ú}{{\'U}}1 {Ñ}{{\~N}}1 {ü}{{\"u}}1 {¿}{{?`}}1 {¡}{{!`}}1 {…}{{...}}3
           {θ}{{$\theta$}}1 {·}{{$\cdot$}}1 {«}{{\guillemotleft}}1 {»}{{\guillemotright}}1 {°}{{$^\circ$}}1}
\usepackage{hyperref}
\hypersetup{colorlinks=true, linkcolor=naranja, urlcolor=naranja}
\usepackage{fancyhdr}
\pagestyle{fancy}
\fancyhf{}
\cfoot{\small \thepage}
\newcommand{\codigo}[1]{\texttt{#1}}
\newcommand{\archivo}[1]{\textbf{\texttt{#1}}}
\newcommand{\aviso}[1]{\par\smallskip\noindent\colorbox{rojo!10}{\parbox{\dimexpr\linewidth-2\fboxsep}{\textbf{\color{rojo}Atención.} #1}}\par\smallskip}
\newcommand{\nota}[1]{\par\smallskip\noindent\colorbox{gris}{\parbox{\dimexpr\linewidth-2\fboxsep}{#1}}\par\smallskip}
\tikzset{
  marca/.style={circle, fill=naranja, text=white, font=\bfseries\footnotesize, inner sep=1.4pt, minimum size=4.6mm, draw=white, line width=0.8pt},
  caja/.style={draw, rounded corners=2pt, fill=gris, align=center, font=\small, minimum height=10mm, inner sep=4pt},
  cajab/.style={caja, fill=white},
  cajan/.style={caja, fill=naranja!14, draw=naranja!70!black},
  paso/.style={draw, rounded corners=2pt, fill=white, align=left, font=\small, text width=3.4cm, minimum height=15mm, inner sep=5pt},
  num/.style={circle, fill=naranja, text=white, font=\bfseries\small, inner sep=1.5pt, minimum size=6mm},
  flecha/.style={-{Stealth[length=2.2mm]}, thick},
  dflecha/.style={{Stealth[length=2.2mm]}-{Stealth[length=2.2mm]}, thick},
  etiq/.style={font=\scriptsize, align=center},
}
\DeclareRobustCommand{\num}[1]{\tikz[baseline=-0.7ex]\node[marca]{#1};}
% \captura{ancho}{archivo}{marcas en coordenadas 0-1 desde la esquina inferior izquierda}
\newcommand{\captura}[3]{%
  \begin{tikzpicture}
  \node[anchor=south west, inner sep=0, draw=black!25] (img) {\includegraphics[width=#1]{#2}};
  \begin{scope}[x={(img.south east)}, y={(img.north west)}]
  #3
  \end{scope}
  \end{tikzpicture}}
% ventana dibujada (consolas y diálogos): \ventana{ancho}{alto}{título}{contenido}
\newcommand{\ventana}[4]{%
\begin{tikzpicture}
\node[draw=black!50, fill=white, minimum width=#1, minimum height=#2, anchor=north west, inner sep=0] (v) at (0,0) {};
\fill[black!8] (v.north west) rectangle ($(v.north east)+(0,-6mm)$);
\node[anchor=west, font=\scriptsize] at ($(v.north west)+(2mm,-3mm)$) {#3};
\node[anchor=east, font=\scriptsize] at ($(v.north east)+(-2mm,-3mm)$) {--\quad$\square$\quad$\times$};
\node[anchor=north west, inner sep=3mm, text width=#1-6mm, font=\scriptsize\ttfamily, align=left] at ($(v.north west)+(0,-6mm)$) {#4};
\end{tikzpicture}}
